\providecommand{\WReach}{\operatorname{WReach}}
\providecommand{\wcol}{\operatorname{wcol}}
\section{Dynamic approximation of distance-\texorpdfstring{$1$}{1} dominating sets}\label{sec:apx}

Here we give a dynamic constant-factor approximation for the domination number in graphs of bounded degeneracy.
We first explain the general weak-reachability argument behind the approximation and only afterwards specialize it to the distance-$1$ representation used by the dynamic implementation. The reason for this is that the general argument applies to an arbitrary distance, but we are able to dynamically maintain the relevant structures only for distance $1$.
Throughout this section, for a graph $G$ we denote by $\gamma(G)$ the minimum size of a (distance-$1$) dominating set in $G$. A \emph{set system} is a family of subsets of some universe $U$, and a \emph{packing} in a set system $\mathcal{F}$ is a subfamily $\mathcal{M}\subseteq \mathcal{F}$ consisting of pairwise disjoint members of $\mathcal{F}$. A packing $\mathcal{M}$ is \emph{maximal} if no superset of $\mathcal{M}$ is a packing; equivalently, every set $F\in \mathcal{F}\setminus \mathcal{M}$ intersects a member of $\mathcal{M}$. The \emph{arity} of a set system $\mathcal{F}$ is the maximum cardinality of a member of $\mathcal{F}$.

\subsection{Combinatorial core}

Fix a graph $G$, an ordering $\sigma$ of $V(G)$, and a radius $r\in\N$.
For $v\in V(G)$, let $\WReach_r[G,\sigma,v]$ be the set of vertices $u\in V(G)$ for which there exists a path $u=p_0,\ldots,p_\ell=v$ of length $\ell\le r$ such that $\sigma(u)\le \sigma(p_i)$ for every $i\in\set{1,\ldots,\ell}$.
Further, let $\wcol_r(G,\sigma)\coloneqq\max_{v\in V(G)} |\WReach_r[G,\sigma,v]|$ and $\mathcal{F}_r(\sigma)\coloneqq\setof{\WReach_r[G,\sigma,v]}{v\in V(G)}$.
The set system $\mathcal{F}_r(\sigma)$ is the general combinatorial source of the approximation: it works for every radius $r$, whereas the dynamic construction below will retain only the case $r=1$.
The next two lemmas formalize the hitting-packing argument showing that weak reachable sets yield a constant-factor approximation of minimum distance-$r$ domination.

\begin{lemma}\label{lem:apx-wreach-domination}
Let $\mathcal{M}\subseteq \mathcal{F}_r(\sigma)$ be a maximal packing in $\mathcal{F}_r(\sigma)$.
Then $D\coloneqq\bigcup_{X\in \mathcal{M}} X$ is a distance-$r$ dominating set in $G$.
\end{lemma}
\begin{proof}
Take any vertex $x\in V(G)$.
If $x\in D$, then $x$ is distance-$r$ dominated by itself.
So suppose that $x\notin D$.
As $\mathcal{M}$ is maximal, the set $\WReach_r[G,\sigma,x]\in \mathcal{F}_r(\sigma)$ intersects some member of $\mathcal{M}$.
Choose $X\in \mathcal{M}$ and $d\in \WReach_r[G,\sigma,x]\cap X$.
By the definition of $\WReach_r[G,\sigma,x]$, there is a $d$-$x$ path of length at most $r$.
Since $d\in X\subseteq D$, this shows that $x$ is distance-$r$ dominated by a vertex of~$D$.
\end{proof}

\begin{lemma}\label{lem:apx-wreach-transversal}
Let $D^*$ be a distance-$r$ dominating set in $G$. Define $T\coloneqq\bigcup_{d\in D^*} \WReach_r[G,\sigma,d]$.
Then $T$ intersects every set in $\mathcal{F}_r(\sigma)$.
\end{lemma}
\begin{proof}
Fix any vertex $v\in V(G)$.
Since $D^*$ is a distance-$r$ dominating set, there exists $d\in D^*$ with $\dist_G(v,d)\le r$.
Let $P$ be a shortest $v$-$d$ path and let $u$ be the vertex of $P$ with minimum position in $\sigma$.
The subpath of $P$ from $u$ to $v$ has length at most $r$, and every internal vertex of this subpath has position at least $\sigma(u)$.
Hence $u\in \WReach_r[G,\sigma,v]$.
By the same argument applied to the subpath of $P$ from $u$ to~$d$, we obtain $u\in \WReach_r[G,\sigma,d]\subseteq T$.
Therefore $\WReach_r[G,\sigma,v]\cap T\neq \emptyset$.
\end{proof}

\begin{lemma}\label{lem:apx-wreach-static}
Let $\mathcal{M}\subseteq \mathcal{F}_r(\sigma)$ be a maximal packing in $\mathcal{F}_r(\sigma)$, $D\coloneqq\bigcup_{X\in \mathcal{M}} X$, and $D^*$ be a minimum distance-$r$ dominating set in $G$.
Then $D$ is a distance-$r$ dominating set in $G$ satisfying \mbox{$|D|\le \wcol_r(G,\sigma)^2\cdot |D^*|$.}
Consequently, whenever $\wcol_r(G,\sigma)\le c$, the set $D$ is a distance-$r$ dominating set of size at most $c^2$ times the optimum.
\end{lemma}
\begin{proof}
By \cref{lem:apx-wreach-domination}, the set $D$ is a distance-$r$ dominating set.

Let $T\coloneqq\bigcup_{d\in D^*} \WReach_r[G,\sigma,d]$.
By \cref{lem:apx-wreach-transversal}, every set in $\mathcal{M}$ intersects $T$.
Since $\mathcal{M}$ is a packing, its sets are pairwise disjoint, hence $|\mathcal{M}|\le |T|$.
Using the definition of $\wcol_r(G,\sigma)$, we infer that
\[|\mathcal{M}|\le |T|\le \sum_{d\in D^*} |\WReach_r[G,\sigma,d]|\le \wcol_r(G,\sigma)\cdot |D^*|.\]
Again by the definition of $\wcol_r(G,\sigma)$, every set in $\mathcal{M}$ has size at most $\wcol_r(G,\sigma)$, so
\[|D| = \sum_{X\in \mathcal{M}} |X| \le \wcol_r(G,\sigma)\cdot |\mathcal{M}| \le \wcol_r(G,\sigma)^2\cdot |D^*|.\qedhere\]
\end{proof}

Thus, \cref{lem:apx-wreach-static} already contains the whole approximation argument.
Note that if $\vec{G}_\sigma$ is the orientation obtained by directing every edge towards the smaller endpoint in $\sigma$, then \[\WReach_1[G,\sigma,v]=\set{v}\cup N^+_{\vec{G}_\sigma}(v)\qquad\textrm{for every vertex }v.\]
In our dynamic data structure, we do not maintain an order $\sigma$ explicitly. Instead, we maintain an orientation of bounded outdegree using the data structure of Brodal and Fagerberg~\cite{BrodalF99}.
Therefore, we now adjust the result of \cref{lem:apx-wreach-static} to the setting of distance-$1$ dominating sets and orientations of bounded outdegree.
The sets $S(v)$ introduced below are not a substitute for the preceding $\WReach_r$ argument; they are exactly its distance-$1$ implementation-level representation.

Let $G$ be a graph and let $\vec{G}$ be an orientation of $G$.
Recall that for every vertex $v\in V(G)$, by $N^+_{\vec{G}}(v)$ we denote the out-neighborhood of $v$ in $\vec{G}$. Define \[S(v)\coloneqq\set{v}\cup N^+_{\vec{G}}(v).\]
We regard the family ${\cal S}\coloneqq \setof{S(v)}{v\in V(G)}$ as a set system over the vertex set $V(G)$, with member per each vertex of $G$.

\begin{lemma}\label{lem:apx-domination}
Let $M\subseteq V(G)$ be such that $\setof{S(v)}{v\in M}$ is a maximal packing in the set system ${\cal S}$.
Then $D\coloneqq\bigcup_{v\in M} S(v)$ is a dominating set in $G$.
\end{lemma}
\begin{proof}
Take any vertex $x\in V(G)$.
If $x\in D$, then $x$ is dominated.
Assume now that $x\notin D$.
By maximality, the set $S(x)$ is not disjoint from at least one set $S(v)$ with $v\in M$.
Choose $v\in M$ and $y\in S(x)\cap S(v)$.
Since $y\in S(v)\subseteq D$, we have $y\neq x$.
Therefore $y\in N^+_{\vec{G}}(x)$, which implies that $xy\in E(G)$.
Hence $x$ is dominated by the vertex $y\in D$.
\end{proof}

\begin{lemma}\label{lem:apx-transversal}
Let $D^*$ be a dominating set in $G$ and define $T\coloneqq\bigcup_{d\in D^*} S(d)$.
Then $T$ intersects every set $S(v)$ with $v\in V(G)$.
\end{lemma}
\begin{proof}
Fix any vertex $v\in V(G)$.
Since $D^*$ is a dominating set, there exists $d\in D^*$ with $d\in N[v]$.
If $d=v$, then $v\in S(v)\cap T$.
Assume therefore that $d\neq v$.
If the edge $vd$ is oriented in $\vec{G}$ from $v$ to $d$, then $d\in S(v)\cap S(d)\subseteq S(v)\cap T$.
Otherwise the edge $vd$ is oriented from $d$ to $v$, and then $v\in S(v)\cap S(d)\subseteq S(v)\cap T$.
Thus $S(v)\cap T\neq \emptyset$.
\end{proof}

\begin{lemma}\label{lem:apx-static}
Assume that $\vec{G}$ has maximum out-degree at most $\Delta$.
Let $M\subseteq V(G)$ be such that the family $\setof{S(v)}{v\in M}$ is a maximal packing in ${\cal S}$.
Then $D\coloneqq\bigcup_{v\in M} S(v)$ is a dominating set in $G$ and satisfies $|D|\le (\Delta+1)^2\cdot \gamma(G)$.
\end{lemma}
\begin{proof}
By \cref{lem:apx-domination}, the set $D$ is a dominating set.

Let $D^*$ be a minimum dominating set in $G$ and let $T\coloneqq\bigcup_{d\in D^*} S(d)$.
By \cref{lem:apx-transversal}, every set $S(v)$ with $v\in M$ intersects $T$.
Since the family $\setof{S(v)}{v\in M}$ is a packing in $\cal S$, these sets are pairwise disjoint, and hence $|M|\le |T|$.
Every set $S(d)$ has size at most $\Delta+1$, so \[|M|\le |T|\le \sum_{d\in D^*} |S(d)|\le (\Delta+1)\cdot |D^*| = (\Delta+1)\cdot \gamma(G).\]
Using again that $|S(v)|\le \Delta+1$ for every $v\in V(G)$, we obtain \[|D| = \sum_{v\in M} |S(v)|\le (\Delta+1)\cdot |M| \le (\Delta+1)^2\cdot \gamma(G).\qedhere\]
\end{proof}

\subsection{Dynamic data structure}

We now propose a dynamic data structure that maintains an approximation domination number of a graph of bounded degeneracy based on the approximation method yielded by \cref{lem:apx-static}. We will exploit the following two results.

\wninline{Zdeduplikowac Brodala}
\begin{theorem}[Brodal--Fagerberg orientation data structure,~\cite{BrodalF99}]\label{thm:BF}
	For every $d\in \N$, there is a deterministic dynamic data structure that for a fully dynamic $n$-vertex graph $G$ whose degeneracy is always bounded by $d$, maintains an orientation $\vec{G}$ with maximum out-degree bounded by $4d$, using $\Oh(n+m)$ words of memory, where $m$ is the current edge count of $G$. It can be initialized on the edgeless $n$-vertex graph in time $\Oh(n)$. It supports edge insertions in amortized time $\Oh(1)$ and edge deletions in amortized time $\Oh(d+\log n)$. Moreover, after each update it can report all edges whose orientation changed during this update, and the amortized number of such edges is $\Oh(d+\log n)$.
\end{theorem}

\begin{theorem}[Assadi--Solomon maximal packing data structure,~\cite{AssadiSolomon21}]\label{thm:AS}
	For every $\Delta\in \N$, there is a randomized dynamic data structure that, against an oblivious adversary, maintains a maximal packing in a fully dynamic labelled set system over a universe $U$ and of arity at most $\Delta$. It supports insertions and deletions of sets in expected amortized time $\Oh(\Delta^2)$, with the same bound holding with high probability, and can be implemented using $\Oh(|U|+\Delta N)$ words of space, where $N$ is the number of maintained labelled sets.
\end{theorem}

We note that Brodal and Fagerberg state their result for graphs of bounded \emph{arboricity} instead of degeneracy, however it is well-known that every $d$-degenerate graph has arboricity at most $d$. 
Also, Assadi and Solomon phrase their result in the language of maximal matchings in hypergraphs, which is equivalent to our terminology of maximal packings in set systems. 

Throughout this subsection, we assume $n\ge 2$ and $d\ge 1$, and we consider a fully dynamic $d$-degenerate graph $G$ on the fixed vertex set $V=\set{1,\ldots,n}$.
Hence, we may invoke the Brodal--Fagerberg orientation theorem (\cref{thm:BF}) on $G$, and thus maintain also an orientation $\vec{G}$ of $G$ with maximum out-degree at most $4d$. We follow the convention that $G_t$ is the graph $G$ at time step $t$, and consequently $\vec{G}_t$ is the maintained orientation of $G_t$.
For every vertex $v\in V$, we let $S_t(v)\coloneqq\set{v}\cup N^+_{\vec{G}_t}(v)$.

\paragraph{Persistent state.}
At time step $t$, the complete persistent state of the data structure consists of the following objects.
\begin{itemize}
\item \emph{Orientation component}: The data structure provided by \cref{thm:BF}, storing the current graph $G_t$ together with its orientation $\vec{G}_t$ and, on every graph update, reporting the updated edge together with the list of all edges reoriented during this update.
\item Set system $\mathcal{S}_t\coloneqq\setof{S_t(v)}{v\in V}$ on the universe $V$, stored explicitly. Every set $S_t(v)\in \mathcal{S}_t$ is labelled with a permanent label $v$. For every $v\in V$, we store the current members of $S_t(v)$ explicitly, and also we store the size of $S_t(v)$. Equivalently, we store the permanent element $v$ together with the current out-neighborhood $N^+_{\vec{G}_t}(v)$ and the cardinality of $\{v\}\cup N^+_{\vec{G}_t}(v)$.
\item \emph{Packing component}: The data structure provided by \cref{thm:AS} for the set system $\mathcal{S}_t$, storing a maximal packing $\mathcal{M}_t$ in $\mathcal{S}_t$ as an iterable list of the labels of the members of $\mathcal{M}_t$. Note that by \cref{lem:apx-static}, the set $D_t\coloneqq \bigcup_{X\in \mathcal{M}_t} X$ is a dominating set in $G_t$ of size at most $(4d+1)^2\cdot \gamma(G_t)$. Therefore, together with $\mathcal{M}_t$ we also store the cardinality of $D_t$. This can be easily updated within the data structure of \cref{thm:AS} upon updates to $\mathcal{M}_t$, as with every set in $\mathcal{S}_t$ we explicitly store also its cardinality.
\end{itemize}
No further persistent objects are needed. In particular, the current dominating set is not stored explicitly: it is represented implicitly by the packing $\mathcal{M}_t$.

\paragraph{Interface.}
As announced in \cref{thm:apx-dynamic}, our data structure supports five methods.
\begin{itemize}
\item $\init{n}$: initialize the data structure on $G_0$ being the edgeless graph on vertex set $V$. 
\item $\add{xy}$: insert the edge $xy$ and update all maintained components.
\item $\remove{xy}$: delete the edge $xy$ and update all maintained components.
\item $\mathtt{SizeDominatingSet()}$: return the cardinality of the current dominating set represented by the packing component.
\item $\mathtt{DominatingSet()}$: returns the current dominating set represented by the packing component.
\end{itemize}
Their guarantees and complexities are summarized in \cref{thm:apx-dynamic}. We now describe the implementation.

\paragraph{Synchronization of the orientation and packing components.}
Whenever a graph update touches an edge $xy$, the set system $\mathcal{S}_t=\{S_t(v)\colon v\in V\}$ needs to be modified only by updating the sets indexed by $x$ and $y$.
More precisely, creating an oriented edge $(x,y)$ inserts $y$ into the $S_t(x)$, deleting an oriented edge $(x,y)$ deletes $y$ from $S_t(x)$, and flipping the orientation of $xy$ from $(x,y)$ to $(y,x)$ is implemented as deleting $y$ from $S_t(x)$ and inserting $x$ into $S_t(y)$.
Thus, every changed oriented edge yields at most two updates to $\mathcal{S}_t$, each also relayed immediately to the packing component.
In the packing component, each modification of a labelled set is implemented by deleting its old version and inserting its updated version.
By \cref{thm:BF}, one update to $G_t$ produces $\Oh(d+\log n)$ updates to $\vec{G}_t$ in the amortized sense: the updated graph edge contributes one update to $\mathcal{S}_t$, and every reorientation performed during the update contributes two more.

\paragraph{Method $\init{n}$} is implemented as follows:
\begin{enumerate}
\item Initialize the orientation component on the edgeless graph on $V=\{1,\ldots,n\}$.
\item For every vertex $v\in V$, create the singleton set $S(v)=\set{v}$ labelled with $v$.
\item Initialize the packing component on the set system consisting of these $n$ labelled sets.
\end{enumerate}
After these steps the orientation is correct, the stored set system is exactly $\mathcal{S}_0=\setof{\set{v}}{v\in V}$, and the packing component stores a maximal packing in $\mathcal{S}_0$.
The initialization time is $\Oh(n)$, because we create $n$ singleton sets.

\paragraph{Method $\add{xy}$} is implemented as follows:
\begin{enumerate}
\item Execute the insertion of $xy$ in the orientation component.
\item Read the newly oriented copy of $xy$ and the list of edges reoriented during this update.
\item For every changed oriented edge (inserted, removed, or reoriented), relay the modification to appropriately update $\mathcal{S}_t$ and the packing component, as described in the Synchronization paragraph above.
\end{enumerate}
Only the sets labelled with endpoints of changed oriented edges are modified, hence after Step~3 every set labelled with $v\in V$ stored in the set system again equals the current set $S_t(v)$, hence the set system is indeed equal to $\mathcal{S}_t$.
Since the packing component is updated after every modification of $\mathcal{S}_t$, at the end it stores a maximal packing in the updated set system.
The amortized number of updates to $\mathcal{S}_t$ is $\Oh(d+\log n)$.
Each of them touches a set of size at most $4d+1$, hence, by \cref{thm:AS}, it costs an expected amortized time of $\Oh((4d+1)^2)=\Oh(d^2)$.
Therefore $\add{xy}$ runs in expected amortized time $\Oh(d^2(d+\log n))$.

\paragraph{Method $\remove{xy}$} is implemented as follows:
\begin{enumerate}
\item Execute the deletion of $xy$ in the orientation component.
\item Read the deleted oriented copy of $xy$ and the list of edges reoriented during this update.
\item For every changed oriented edge, relay the modification to appropriately update $\mathcal{S}_t$ and the packing component, as described in the Synchronization paragraph above.
\end{enumerate}
The same argument as in the $\add{xy}$ method shows that after Step~3 the stored set system is exactly $\mathcal{S}_t$, and the packing component stores a maximal packing in it.
Again there are amortized $\Oh(d+\log n)$ updates to $\mathcal{S}_t$, each affecting a set of size at most $4d+1$, so $\remove{xy}$ also runs in expected amortized time $\Oh(d^2(d+\log n))$.

\paragraph{Method $\mathtt{SizeDominatingSet()}$} is implemented by just returning the cardinality of $D_t=\bigcup_{X \in \mathcal{M}_t} X$, stored explicitly within the packing component.


\paragraph{Method $\mathtt{DominatingSet()}$} is implemented as follows.
We traverse the set system $\mathcal{M}_t$ (stored as an iterable list of labels of the members of $\mathcal{M}_t$) and we output the union $D_t=\bigcup_{X\in \mathcal{M}_t} X$.
No duplicate elimination is needed, because the members of $\mathcal{M}_t$ are pairwise disjoint.
%By \cref{lem:apx-static} applied with $\Delta\coloneqq 4d$, the returned set $D_t$ is a dominating set in $G_t$ and satisfies $|D_t|\le (4d+1)^2\cdot \gamma(G_t)$.
Since the query outputs exactly the vertices contained in the members of $\mathcal{M}_t$, its running time is $\Oh(|D_t|)$.

\paragraph{Memory usage.}
Let $m_t\coloneqq|E(G_t)|$. By \cref{thm:BF},
the orientation component uses $\Oh(n+m_t)$ words.
The explicitly stored set system $\mathcal{S}_t$ also uses $\Oh(n+m_t)$ words, because the total sum of cardinalities of the members of $\mathcal{S}_t$ is exactly $n+m_t$. %every member of $\mathcal{S}_t$ contributes its permanent vertex label once, accounting for $n$ stored incidences, and every graph edge contributes exactly one additional stored incidence, namely to the hyperedge of its tail in $\vec{G}_t$.
The packing component is maintained over the universe $V$, with $n$ labelled sets of arity at most $4d+1$, so by \cref{thm:AS}, it uses $\Oh(n+(4d+1)n)=\Oh(dn)$ words.
The dominating set returned by $\mathtt{DominatingSet()}$ is not stored explicitly, so it contributes no additional space.
Hence the whole data structure occupies $\Oh(n+m_t)+\Oh(n+m_t)+\Oh(dn)=\Oh(dn+m_t)$ words.

\michalin{The theorem below should go to intro. We should explain there also what is an oblivious adversary.}

The description provided above amounts to the proof of the following result, which summarizes our data structure.

\begin{theorem}\label{thm:apx-dynamic}
For every $d\in \N$ there is a randomized dynamic data structure that, against an oblivious adversary, maintains a fully dynamic $d$-degenerate graph $G$ on $n$ vertices together with a dominating set $D$ in $G$ satisfying $|D|\leq (4d+1)^2\cdot \gamma(G)$. The access to $D$ is provided by queries that report its size in time $\Oh(1)$ and enumerate all its vertices in time $\Oh(|D|)$. The data structure can be initialized on an edgeless $n$-vertex graph in time $\Oh(n)$, uses $\Oh(dn+m)$ words of space where $m$ denotes the current number of edges, and processes each edge insertion and deletion in expected amortized time $\Oh(d^2(d+\log n))$.
\end{theorem}
\begin{comment}
\begin{proof}
The description above already specifies the whole persistent state, all supported methods, and their local invariants.
Initialization takes $\Oh(n)$ time by the analysis of $\init{G_0}$.
At every later moment $t$, the query procedure returns $D_t\coloneqq\bigcup_{v\in M_t} S_t(v)$, where $M_t$ is the label set of the maintained maximal matching.
Because $\vec{G}_t$ has maximum out-degree at most $4d$, \cref{lem:apx-static} yields that $D_t$ is a dominating set in $G_t$ and satisfies $|D_t|\le (4d+1)^2\cdot \gamma(G_t)$.
The matched hyperedges are pairwise disjoint, so reporting their union takes time $\Oh(|D_t|)$.

For updates, the orientation algorithm is deterministic, hence the sequence of hypergraph updates depends only on the graph update sequence and not on the random choices of the matching component.
Therefore the oblivious-adversary guarantee of Assadi--Solomon applies to the whole execution.
The analyses of $\add{xy}$ and $\remove{xy}$ show that each graph update performs amortized $\Oh(d+\log n)$ hypergraph updates, each of expected amortized cost $\Oh((4d+1)^2)$.
Thus every edge insertion and deletion runs in expected amortized time $\Oh(d^2(d+\log n))$.

Finally, the memory analysis above gives the space bound $\Oh(dn+m_t)$ at time $t$, which is exactly the claimed $\Oh(dn+m)$ bound.
\end{proof}
\end{comment}
